\lab{5}{Закон Ципфа}
\task
Построить распределение частот терминов в логарифмической шкале,
наложить закон Ципфа и объяснить расхождения. Необязательная часть -
подбор параметров закона Мандельброта.

\method
Для каждого уникального токена считаются число вхождений \(cf\) и число
документов \(df\). По убыванию \(cf\) назначаются ранги \(r=1,\ldots,V\);
равные частоты упорядочиваются лексикографически для воспроизводимости.
В этой работе морфологическая нормализация не включается, чтобы
отделить свойства токенизации от стемминга.

Модель Ципфа имеет вид \(f(r)=C/r\). Константа выбирается из сохранения
общей массы частот:
\[
H_V=\sum_{r=1}^{V}\frac1r,\qquad C=\frac{T}{H_V},
\qquad\sum_{r=1}^{V}\frac{C}{r}=T.
\]
Это не подгонка по самому частому слову. Дополнительно методом наименьших
квадратов оценивается наклон зависимости \(\ln f\) от \(\ln r\)
по всем термам. В идеальной модели он равен \(-1\).
Для отображения можно проредить ранги; расчёт регрессии выполняется
до прореживания. Гапакс - термин с одним вхождением во всём наборе.

\journal
Потоковый токенизатор дополнен частотами и множеством термов текущего
документа, чтобы \(df\) увеличивалась лишь один раз на документ.
Скрипт \code{5_zipf/tools/zipf.py} проверяет уникальность строк,
\(cf\ge df\ge1\), сортирует термы, считает модель и регрессию.
Статистика вычислена для первых 10\,000 документов
корпуса. Поскольку регрессия и сортировка требуют полного словаря
исследуемого набора, для этого опыта выбран ограниченный префикс.
Полученные частоты и наклон характеризуют именно эти 10\,000 документов.

\testing
Для искусственного словаря проверяются ранги, сумма частот и значение
константы. Пустой словарь, дубликат терма и \(cf<df\) должны отвергаться.
При совпадении частот порядок должен оставаться стабильным. Для
одного терма наклон условно равен нулю, поскольку дисперсия логарифмов
ранга равна нулю. При прореживании необходимо сохранять первый и
последний ранги. Перечисленные проверки задают план тестирования
частотного анализа и построения графика.

\results
В эксперименте: 10\,000 документов, 10\,683\,808 вхождений,
401\,968 уникальных токенов, 266\,496 гапакс
(66.30\% словаря).
Константа $C=792488.31$, наклон регрессии $-1.0655$.
Это характеристики выборки без стемминга, а не всего корпуса.
Самые частые токены: \code{the} (544\,244), \code{to} (230\,087), \code{a} (192\,740), \code{is} (175\,251), \code{of} (170\,951).

\begin{figure}[H]\centering
\begin{tikzpicture}
\begin{loglogaxis}[width=0.95\textwidth,height=7cm,
xlabel={Ранг термина},ylabel={Число вхождений},grid=both,
legend style={at={(0.02,0.02)},anchor=south west,font=\small}]
\addplot[blue,thick] table[x=rank,y=count]{figures/zipf.dat};
\addlegendentry{Наблюдаемые частоты}
\addplot[red,dashed,thick] table[x=rank,y=zipf]{figures/zipf.dat};
\addlegendentry{Модель $C/r$}
\end{loglogaxis}
\end{tikzpicture}
\caption{Распределение для первых 10\,000 документов. Все точки получены из фактических частот этой выборки.}
\end{figure}
Частые слова образуют начало распределения, технические имена и
уникальные обозначения - длинный хвост. При малых целых частотах
видны ступени. Отклонение от прямой объясняется смесью прозы и кода,
ограниченностью выборки, повторяемостью документации и отсутствием
фильтра стоп-слов. Линейная регрессия в логарифмах является описанием
формы кривой, а не статистическим доказательством степенного закона.

\conclusions
Распределение имеет длинный хвост редких терминов, поэтому размер словаря
может расти значительно быстрее, чем подсказывает число частых слов.
Модель Ципфа полезна как приближение и объясняет неравномерную длину
списков постингов. Опыт ограничен префиксом корпуса с документацией;
для вывода о всём корпусе нужен полный пересчёт или стратифицированная
выборка обсуждений и документации. Закон Мандельброта не подбирался.
